\chapter{The Volume Problem and Isoperimetry}
\label{ch:volume}

The volume of a convex body is a $\sharpP$ hard problem for which the full theory exists: a proof that deterministic approximation is impossible in the oracle model, a randomized polynomial time algorithm, and thirty years of improvement in the exponent driven by progress on a single inequality. It is the model for what a mixing proof for the queens would look like, and this chapter gives the theory in more detail than Chapter~\ref{ch:labs} could.

\section{The lower bound}

A convex body $K \subseteq \mathbb{R}^n$ is given by a membership oracle, together with bounding balls $B(0, r) \subseteq K \subseteq B(0, R)$. Elekes showed in 1986 that any deterministic algorithm making polynomially many oracle calls cannot approximate $\mathrm{vol}(K)$ to within a factor better than $2^{n}$, and B\'ar\'any and F\"uredi sharpened this in 1987 to $(cn / \log n)^{n/2}$. The argument is that the convex hull of polynomially many query points has volume exponentially smaller than the ball, so that the algorithm cannot distinguish the ball from the hull of its own queries. The lower bound is unconditional and does not depend on $\Pc \neq \NP$; it is a statement about oracle algorithms, and it says that randomness is provably necessary for approximating volume in the oracle model.

\section{Dyer, Frieze and Kannan}

The randomized algorithm of 1991 proceeds by a ladder of bodies $K_i = K \cap B(0, r 2^{i/n})$ for $i = 0, \dots, m$ with $m = n \log_2(R/r)$, so that $K_0 = B(0, r)$ has known volume, $K_m = K$, and consecutive ratios $\mathrm{vol}(K_i) / \mathrm{vol}(K_{i+1})$ are at least $\tfrac12$ by the scaling of volume in dimension $n$. Each ratio is estimated by sampling from $K_{i+1}$ and counting the fraction in $K_i$. The sampler is a random walk on a fine grid inside $K_{i+1}$, and the proof that it mixes in polynomial time goes through an isoperimetric inequality: for a convex body $K$ and a partition $K = S_1 \cup S_2 \cup S_3$ with $S_3$ separating $S_1$ from $S_2$, the volume of $S_3$ is at least a constant times the distance between $S_1$ and $S_2$ divided by the diameter, times the smaller of the volumes of $S_1$ and $S_2$. This is a geometric statement that convex bodies have no narrow bottlenecks, and it converts to a conductance bound for the walk by the Jerrum Sinclair argument.

The original running time was $O^*(n^{23})$. Successive improvements came from better walks, the ball walk and the hit and run walk; better isoperimetric inequalities, in particular the localisation lemma of Lov\'asz and Simonovits which reduces an $n$ dimensional isoperimetric statement to a one dimensional one; and better ladders, including the annealing of Lov\'asz and Vempala, in which the body is fixed and an exponential weight is annealed, and the Gaussian cooling of Cousins and Vempala~\cite{CousinsVempala} (2018). Lov\'asz and Vempala reached $O^*(n^4)$ in 2006 and the current record is $O^*(n^3)$, up to factors governed by the KLS constant.

\section{The log concave extension}

Applegate and Kannan in 1991, and Frieze, Kannan and Polson in 1994, extended the analysis from indicators of convex bodies to log concave functions, those with $f = e^{-\phi}$ for convex $\phi$. The point is that a log concave density is the natural generalisation of a uniform density on a convex set, and that every Bayesian posterior with log concave likelihood and prior is of this form. Applegate and Kannan used a grid walk with Metropolis acceptance and proved polynomial mixing from a warm start by a discretised isoperimetric inequality. Frieze, Kannan and Polson analysed the Metropolis ball walk directly on the continuous space, proving a conductance bound from the isoperimetry of log concave measures, and obtained a polynomial mixing time with explicit dependence on the dimension, the diameter of the support and the smoothness of $\phi$. Their paper established that posterior sampling for log concave models is a problem with a polynomial time guarantee.

The vertical ladder for the evidence of a log concave model inherits this guarantee: each level set of a log concave likelihood under a log concave prior is a convex body weighted by a log concave density, and the Frieze Kannan Polson walk mixes on it in polynomial time. The Bayesian evidence computation for such models is therefore an FPRAS, and this is the case, discussed in Chapter~\ref{ch:bayes}, where the program succeeds in full.

\section{The KLS conjecture}

The isoperimetric constant that governs all these walks is the Cheeger constant of the body or density: the infimum over partitions of the boundary measure divided by the smaller volume. Kannan, Lov\'asz and Simonovits conjectured in 1995 that for an isotropic log concave measure, one with identity covariance, the Cheeger constant is bounded below by a universal constant independent of dimension. If true, this gives the best possible mixing for the ball walk, of order $n^2$ steps from a warm start, and the best possible volume algorithm. The conjecture resisted attack for twenty five years; Eldan reduced it to a question about the thin shell constant in 2013, Lee and Vempala proved a bound of $n^{1/4}$ in 2017, Chen proved $n^{o(1)}$ in 2021 by a stochastic localisation argument, and Klartag and Lehec proved a polylogarithmic bound in 2022, with Klartag~\cite{Klartag2023} reaching $\sqrt{\log n}$ in 2023. The conjecture is now known up to a factor polylogarithmic in $n$, which for the volume algorithm is as good as true.

The progress on KLS is the model for what the program needs. The question for convex bodies was whether the level sets have bottlenecks; the answer, pursued over three decades, is that they do not, up to a factor that is at most polylogarithmic in the dimension (the conjecture of a dimension free constant remains open). The question for the queens is the same, with the level set a subset of $S_N$ under swap adjacency rather than a convex body under Euclidean adjacency. There is no convexity, and the KLS methods, which are all about convexity, do not apply directly. What the giant component computations of Chapter~\ref{ch:open} suggest is that the queens level sets have, for unconditioned boards, an isoperimetric structure of their own: the solutions are densely bridged by near solutions, and the bottleneck, if there is one, is the stragglers. A proof of a Cheeger bound for the giant component would be the analogue of Dyer, Frieze and Kannan for the queens, and we do not have it.

\section{What distinguishes volume from the queens}

Three things. Volume has the lower bound of B\'ar\'any and F\"uredi, which says deterministic methods cannot work; the queens have no such theorem, and if $\Pc = \RP$ a deterministic method exists in principle. Volume has convexity, which gives isoperimetry for free; the queens have a combinatorial structure whose isoperimetry is unknown. And volume has a decision problem in $\Pc$, since membership is given; the queens completion problem is $\NP$ complete. The first two differences are technical and the third is the whole matter: by the main equivalence of Chapter~\ref{ch:theorems}, an isoperimetric bound for the queens level sets on all boards would imply $\NP = \RP$, and so the method that worked for volume cannot work for the queens unless the world is not as expected. The method can work, and we think does, for unconditioned boards, which are the easy case, and the program's task is to say exactly where the method stops.
